% AI-produced by GPT-6 Astra (OpenAI, ChatGPT Pro mode).
% Reconstructed 2026-11-20 from maximal_locus2's complete three-part A151
% authored mathematical text and the subsequently audited radial-collapse
% addendum. Historical TeX bytes were lost; this is a new source identity.
\Needspace{10\baselineskip}
\subsection{A.151. A moving angular tent and the limits of homeomorphic clocks}

\begin{saresult}{A.151}{Proved}{A based semialgebraic family with no homeomorphic Hausdorff-Lipschitz reparameterization}
Work over the real field. Put
\[
 C=\partial[-1,1]^2,\quad B=C^d\ (d\ge1),\quad
 X=B\times[-1,1],\quad P=[-1,1]^2.
\]
There is an explicit jointly continuous semialgebraic family
$U_p:X\to X$ with the following properties.
\begin{enumerate}
\item Every $U_p$ and its inverse are spatially Lipschitz with one
uniform Euclidean constant strictly less than $2$. They are the
identity on both outer half-interval collars in the last coordinate.
\item The family is the identity at $p=0$ and on $\partial P$, and
has an explicit semialgebraic nullhomotopy relative to those parameters
and the spatial boundary.
\item No homeomorphism $h:P\to P$, even a nondefinable TOP
homeomorphism, makes $p\mapsto\Gamma U_{h(p)}$ Lipschitz in the
Euclidean Hausdorff metric.
\item For no such $h$ can one write $U_{h(p)}=s_pr_p^{-1}$ with
two families of homeomorphisms of $X$ that are both uniformly
forward parameter-Lipschitz. This impossibility does not require
the factors to be definable or to preserve any marks.
\item There is a noninjective Lipschitz semialgebraic
$g:P\to P$, equal to the identity on $\partial P$ and at $0$,
such that $U\circ g$ is uniformly forward parameter-Lipschitz.
The straight homotopy from the identity to $g$ gives the allowed
based homotopy from $U$ to this representative.
\end{enumerate}
In particular the example is available on the literal PL six-torus
$B=C^6$. It obstructs a homeomorphic reparameterization or an exact
factorization of the given representative; it is not an obstruction
to the homotopy regularization established later.
\end{saresult}

\begin{proof}
\textbf{The tent and the homeomorphisms.}
For $p\ne0$ define
\[
 r=\|p\|_\infty,\qquad \theta=p/r\in C,\qquad
 a=r(1-r),\qquad
 f_p(z)=\max\{0,a-\|z-\theta\|_\infty\}\quad(z\in C),
 \tag{151.1}
\]
and put $f_0=0$. These are rational semialgebraic formulas, using
division only on $r>0$. The identity
\[
           \frac14-r(1-r)=(r-\tfrac12)^2
 \tag{151.2}
\]
shows $0\le f_p\le a\le1/4$. The bound by $a$ proves continuity
at $p=0$ and at $\partial P$, where $a=0$. Elsewhere continuity
follows directly from the formula. Every spatial slice $f_p$ is
$1$-Lipschitz in the maximum norm and hence also in the Euclidean
norm of the ambient plane.

Write $b=(z_1,\ldots,z_d)\in B$, and set
\[
 \eta(y)=\max\{0,1-2|y|\},\qquad \varepsilon=\tfrac14,\qquad
 U_p(b,y)=\bigl(b,y+\varepsilon f_p(z_1)\eta(y)\bigr).
 \tag{151.3}
\]
For fixed $b,p$, put $c=\varepsilon f_p(z_1)$, so
$0\le c\le1/16$. The last-coordinate map is
\[
 T_c(y)=
 \begin{cases}
 y,&-1\le y\le-\tfrac12,\\
 (1+2c)y+c,&-\tfrac12\le y\le0,\\
 (1-2c)y+c,&0\le y\le\tfrac12,\\
 y,&\tfrac12\le y\le1.
 \end{cases}
 \tag{151.4}
\]
The common join values are $-1/2,c,1/2$, and all slopes belong to
$[7/8,9/8]$. Therefore $T_c$ is an increasing homeomorphism.
Its inverse is the identity outside $[-1/2,1/2]$, and on the
two remaining intervals it is
\[
 T_c^{-1}(Y)=
 \begin{cases}
 (Y-c)/(1+2c),&-\tfrac12\le Y\le c,\\
 (Y-c)/(1-2c),&c\le Y\le\tfrac12.
 \end{cases}
 \tag{151.5}
\]
Positive denominators and the matching join values show that
this inverse is jointly continuous and semialgebraic. Thus
(151.3) is a genuine joint family of homeomorphisms.

We give uniform spatial estimates, independent of $d$.
For two input points let $D=\|b-b'\|_2$ and $E=|y-y'|$.
The last-coordinate difference is at most
$D/4+9E/8$. Consequently
\[
 \|U_p(b,y)-U_p(b',y')\|_2
 \le \tfrac54D+\tfrac98E
 \le\tfrac{5\sqrt2}{4}\sqrt{D^2+E^2}<2\sqrt{D^2+E^2}.
 \tag{151.6}
\]
For the inverse, monotonicity with slope at least $7/8$ gives
\[
 |y-y'|\le\tfrac87|Y-Y'|+\tfrac27\|b-b'\|_2.
\]
The inverse Euclidean Lipschitz constant is therefore at most
$9\sqrt2/7<2$. The strict inequalities follow from
$25/8<4$ and $162/49<4$.

\textbf{The local angular lower bound.}
Let $p,q$ be nonzero interior parameters. Write their amplitudes as
$a,a'$ and their directions as $\theta,\theta'$. If
\[
             \delta=\|\theta-\theta'\|_\infty
                            \le\min\{a,a'\},
\]
evaluation of the two tents at $\theta$ and $\theta'$ gives
\[
 \|f_p-f_q\|_\infty
       \ge\max\{|a-a'+\delta|,\ |a-a'-\delta|\}
       =|a-a'|+\delta\ge\delta.
 \tag{151.7}
\]
The condition on $\delta$ ensures that both subtracted tent values
are on their nonzero affine branches. The equality is the elementary
two-sign identity for real numbers and $\delta\ge0$.

For maps of $X$ define
$d_\infty(V,W)=\sup_{x\in X}\|V(x)-W(x)\|_2$.
Since $\eta(0)=1$ and every $z_1\in C$ occurs in $B$, we have
the exact equality
\[
           d_\infty(U_p,U_q)
                     =\varepsilon\|f_p-f_q\|_\infty.
 \tag{151.8}
\]
Write $d_H$ for the Euclidean Hausdorff distance between their
compact graphs in $X\times X$. Matching the same input point gives
$d_H(\Gamma U_p,\Gamma U_q)\le d_\infty(U_p,U_q)$.
In the other direction, for any $x$ choose a closest point
$(x',U_q(x'))$ in the other graph. The spatial Lipschitz bound $2$
gives
\[
 \begin{split}
 \|U_p(x)-U_q(x)\|_2
 &\le\|U_p(x)-U_q(x')\|_2+
                         2\|x'-x\|_2\\
 &\le3d_H(\Gamma U_p,\Gamma U_q).
 \end{split}
 \tag{151.9}
\]
Thus (151.7)--(151.9) imply
\[
 \|\theta-\theta'\|_\infty
       \le12d_H(\Gamma U_p,\Gamma U_q)
 \quad\text{if }\|\theta-\theta'\|_\infty\le\min\{a,a'\}.
 \tag{151.10}
\]
The graph metric and the forward supremum metric have been kept
distinct; (151.9) uses the uniform spatial bound.

\textbf{An arbitrary parameter homeomorphism cannot repair the metric.}
Suppose, towards a contradiction, that
$p\mapsto\Gamma U_{h(p)}$ is $L$-Lipschitz for a homeomorphism
$h:P\to P$. Put $q=h^{-1}(0)$. A homeomorphism of the closed square
preserves its boundary: invariance of domain applied to the
restrictions of $h$ and $h^{-1}$ to interiors proves that both
send interior points to interior points. Thus $q\in\operatorname{int}P$.

Choose $\rho>0$ so that
\[
 Q_\rho=q+[-\rho,\rho]^2\subset\operatorname{int}P,\qquad
                    \rho<\frac{1}{100(1+L)}.
 \tag{151.11}
\]
On the boundary loop of this small square define
$\alpha(v)=\theta(h(v))$. The image $h(\partial Q_\rho)$ avoids both $0$ and
$\partial P$. Therefore the continuous amplitude on that loop
has a positive minimum $\nu$. Uniform continuity of $\alpha$
allows every boundary arc to be partitioned into finitely many
successive pieces whose angular endpoint distance is at most
$\nu$. Retain the four corners in such a partition.
For every successive pair (151.10) applies, and the hypothesized
Hausdorff estimate bounds its angular distance by
$12L$ times the Euclidean distance between the domain points.
Summing along an arc, whose length is at most the square perimeter
$8\rho$, gives
\[
                   \operatorname{diam}_\infty\alpha(\partial Q_\rho)
                                      \le96L\rho.
 \tag{151.12}
\]
No regularity of $h$ and no length theorem for its image is used:
only finite subdivision of the original square loop is involved.

On the other hand this angular image is all of $C$.
Indeed $K=h(Q_\rho)$ is compact and contains $0$ in its interior,
and $\partial K=h(\partial Q_\rho)$. For each $u\in C$, the ray
$\{tu:t\ge0\}$ has a last point $t_*u$ in $K$; here $t_*>0$.
The last point belongs to $\partial K$, and its maximum-norm
direction is $u$. This proves $\alpha(\partial Q_\rho)=C$.
Its maximum-norm diameter is $2$, contradicting
\[
      2\le96L\rho<
             \frac{96L}{100(1+L)}<\frac{24}{25}<2.
 \tag{151.13}
\]
For $L=0$ the contradiction is already immediate from (151.12).
This proves the assertion for every TOP parameter homeomorphism.

\textbf{Exact factorization would give the forbidden graph bound.}
Suppose $U_{h(p)}=s_pr_p^{-1}$, with forward parameter bounds
$L_r,L_s$ for two families of homeomorphisms. Its graph is
\[
      \Gamma U_{h(p)}
            =\{(r_p(x),s_p(x)):x\in X\}.
\]
Match the same $x$ at two parameter values, in both directions.
This gives
\[
 d_H(\Gamma U_{h(p)},\Gamma U_{h(q)})
                 \le\sqrt{L_r^2+L_s^2}\,\|p-q\|_2,
 \tag{151.14}
\]
which has just been excluded. The reasoning does not use
definability or any marking conditions on the factors.

\textbf{The family nevertheless has an explicit relative nullhomotopy.}
For $0\le\lambda\le1$ replace $\varepsilon$ in (151.3) by
$(1-\lambda)\varepsilon$. The same four slopes still lie in
$[7/8,9/8]$. This is a semialgebraic homotopy through
homeomorphisms from $U$ to the identity, fixed at $p=0$,
on all of $\partial P$, and on the spatial outer collars.
Rescaling $[-1,1]$ affinely to $I$ makes it a based
bottom-fixed concordance cube in the stipulated convention.

\textbf{A noninjective radial clock gives a Lipschitz representative.}
Fix $0<r_0<1$, for example $r_0=1/2$. Define
\[
 R(r)=\max\{0,(r-r_0)/(1-r_0)\},\qquad
 g(p)=
 \begin{cases}
 0,&\|p\|_\infty\le r_0,\\
 R(\|p\|_\infty)\,p/\|p\|_\infty,
                         &\|p\|_\infty\ge r_0 .
 \end{cases}
 \tag{151.15}
\]
The two formulas agree at the interface. This is a continuous
semialgebraic map, collapses the central square to $0$, and fixes
$\partial P$ and $0$.

Here is a uniform estimate that also verifies its Lipschitz property.
For parameters of radii at least $r_0$, put
$\Delta=\|p-q\|_\infty$. Then
\[
 |R(r)-R(r')|\le\Delta/(1-r_0),\qquad
       \|\theta(p)-\theta(q)\|_\infty\le2\Delta/r_0.
 \tag{151.16}
\]
The second inequality follows by subtracting $p/r-q/r'$ and
using $|r-r'|\le\Delta$. Since $0\le R\le1$, these inequalities
make $g$ Lipschitz with constant
$K_0=(1-r_0)^{-1}+2/r_0$ in the maximum norm.
If one radius is below $r_0$, its image is zero and
$R(r)\le|r-r'|/(1-r_0)$ gives the same bound.

The amplitude function $a(u)=u(1-u)$ is $1$-Lipschitz on $I$,
since $|a(u)-a(v)|=|u-v|\,|1-u-v|\le|u-v|$.
Directions in the noncollapsed region are unchanged by $g$.
The tent formula therefore gives
\[
             \|f_{g(p)}-f_{g(q)}\|_\infty
                                     \le K_0\|p-q\|_\infty.
 \tag{151.17}
\]
For a pair straddling the collapsed region use
$0\le f_{g(p)}\le R(r)\le
 \|p-q\|_\infty/(1-r_0)$; inside it both functions vanish.
Combining with (151.8), the forward parameter bound of
$U\circ g$ is $\varepsilon K_0$.
For $r_0=1/2$ this is $3/2$, also a valid Euclidean
parameter bound. The homotopy
$g_t(p)=(1-t)p+tg(p)$ stays in $P$, fixes $\partial P$ and $0$,
and yields the required based homotopy $U\circ g_t$.
For every $t<1$ this parameter map is itself a homeomorphism:
its radial function $(1-t)r+tR(r)$ is strictly increasing from
$0$ to $1$. Only the final map collapses the central square.
The factors $\operatorname{id}$ and $U\circ g$ now satisfy
the forward parameter condition. Thus this example separates
homeomorphic exact reparameterization from the permitted homotopy
regularization.

\textbf{Relation to round circles.}
The main construction already has a literal polyhedral torus base.
If one wants a round-circle model, the maps
\[
 R_1:S^1\to C,\quad R_1(z)=z/\|z\|_\infty,\qquad
 R_1^{-1}(u)=u/\|u\|_2
 \tag{151.18}
\]
are mutually inverse semialgebraic bi-Lipschitz homeomorphisms.
Their constants can be bounded by $\sqrt2+2<4$ and $2$,
respectively, using denominators $\|z\|_\infty\ge1/\sqrt2$
and $\|u\|_2\ge1$. Their finite products, times the interval
identity, remain bi-Lipschitz. For any such spatial conjugacy $J$,
\[
        \Gamma(JVJ^{-1})=(J\times J)(\Gamma V).
\]
Both Lipschitz constants compare the two graph Hausdorff metrics,
so the impossibility of a parameter homeomorphism is preserved.
No equality between a Euclidean-distance tent and the conjugate
of the maximum-norm tent is asserted.
\end{proof}

\textbf{Finite arithmetic certificate and its limits.}
The exact program
\path{agents/singularities_duality/code/check_A151_exact.js}
uses normalized BigInt rational arithmetic and symbolic univariate
polynomials. It checks the amplitude identity, four piece joins,
all endpoint slope bounds, positive inverse denominators, the two
spatial constants, inverse coefficients, the angular and square-loop
constants, and the elementary signed comparisons, in 26 tests.
Its recorded output is
\path{agents/singularities_duality/logs/check_A151_exact.json}.
These are exact finite algebra checks, without floating-point
sampling. Compactness, invariance of the square boundary, existence
of nearest graph points, subdivision of the square loop, and the
ray last-exit argument are the mathematical proofs above, not
claims made by the program. The radial-collapse estimate is also
proved explicitly above. The family has the displayed nullhomotopy,
so this finite example represents no nonzero concordance class.

